
% The environment used here (theappendices) is a wrapper for the basic appendices environment which changes the appearance of the title page and the structure and appearance of the appendices in the table of contents and PDF bookmarks. The original functionality can be restored by simply removing the 'the' from the \begin{} and \end{} statements below.

\begin{appendices}



\section{Transfer Matrix Method python code}
\begin{lstlisting}[language=Python, breaklines=true, breakatwhitespace=true]

# TMM for 9-layers - MJ Hurben Dec 2024
# Modified for 1550 nm and 780 nm optimized SiO2/Si ARC - GCC 2026
import numpy as np
import matplotlib.pyplot as plt
d2r = np.pi / 180  # Degrees to radians factor
na=1               # Index of air
n_layers=3         # Number of layers

# Define matrices of refractive index n and thickness d, modify functions and arrays as needed

# try to use complex refractive index n(wl)+k(wl)j for formula?

def ns(wl): # Define ns as a function of wl based on index of substrate material Si from 430 nm
    return 4.0333E-12*wl**4 - 1.7984E-08*wl**3 + 2.9993E-05*wl**2 - 2.2437E-02*wl + 9.9702
            
#    return -8.974E-15*wl**5 + 5.299E-11*wl**4 - 1.228E-07*wl**3 + 1.398E-04*wl**2 - 7.857E-02*wl + 2.103E+01 # In52Ga48As

def n1(wl): # Define n1 as a function of wl based on refractive index of 1st layer material SiO2
    return np.sqrt(20201856809.3-22607182010.85*wl**2/ (wl**2-0.1063836985797) +2405325203.746*wl**2/ (wl**2-0.9998803009815))

def n2(wl):  # Define n2 as a function of wl a-Si from 430 nm
    return np.sqrt(14.84529+(7067090000000)*wl**2/ (wl**2-(-2124490000000000000)) +(-10.41841)*wl**2/ (wl**2-(-190469.027)))

def refractive_index_array(wl): # Add corresponding variables and links to functions for each layer
    n_s = ns(wl)
    n1_val = n1(wl)
    n2_val = n2(wl)
    return np.array([na, na, na, na, na, na, n1_val, n2_val, n1_val, n_s])   # If < 9 layers, use na as index of 'missing' layers


d=np.array([1,1,1,1,1,1,206,14,53,1000]) # d-Matrix values are in nm; initial air layers and final glass layer arbitrary

def PS(nva,dva,wl): # Phase Shift calculation

    def M(nA,nB): # Transfer matrix for interface
        r=(nA-nB)/(nA+nB)
        t=1+r
        d=1/t
        offd=r/t    
        m=np.array([[d,offd],[offd,d]])
        return m

    def P(L,n):  # Transfer matrix for propagation
        L=L*n*2*np.pi
        v=complex(np.cos(L),-np.sin(L))
        p=np.array([[v,0],[0,v.conjugate()]])
        return p
   
    W=M(nva[8],nva[9]) # Interface of last layer and glass
    for i in range(8,-1,-1):
        W=np.matmul(P(dva[i]/wl,nva[i]),W)
        if i>0:
            W=np.matmul(M(nva[i-1],nva[i]),W)
        else:
            W=np.matmul(M(na,nva[0]),W)  # Interface of glass and first layer
        r=W[1,0]/W[0,0]    
    return r


d_sweep=np.array([1,1,1,1,1,1,206,14,53,1000])   # Matrix for thickness sweep
d_optimal=np.array([1,1,1,1,1,1,206,14,53,1000])

for layer in range(int(9-n_layers),9):   # Rough sweep (100 nm range 5 nm step) for best thickness
    original_val=d_sweep[layer]
    base_reflectance= abs( PS( refractive_index_array(1550), d_sweep,1550))**2 + abs( PS( refractive_index_array(780) ,d_sweep,780)) **2
    for sweep_size in range(-50,55,5):
        d_sweep[layer]=original_val+sweep_size
        r_init_1550= PS(refractive_index_array(1550),d_sweep,1550)
        r_init_780= PS(refractive_index_array(780),d_sweep,780)
        if (abs(r_init_1550)**2+abs(r_init_780)**2) <base_reflectance:   # try reducing sum of base_reflectance at 1550 and 780? 
            d_optimal[layer]=d_sweep[layer]
            base_reflectance = (abs(r_init_1550)**2+abs(r_init_780)**2)
    d_sweep[layer]=d_optimal[layer]

for layer in range(int(9-n_layers),9):   # Fine sweep (10 nm range 1 nm step) for best thickness
    original_val=d_sweep[layer]
    base_reflectance = abs( PS( refractive_index_array(1550), d_sweep,1550)) **2+ abs( PS( refractive_index_array(780), d_sweep,780)) **2
    for sweep_size in range(-5,6):
        d_sweep[layer]=original_val+sweep_size
        r_init_1550 = PS(refractive_index_array(1550),d_sweep,1550)
        r_init_780 = PS(refractive_index_array(780),d_sweep,780)
        if (abs(r_init_1550)**2+abs(r_init_780)**2) <base_reflectance:   # try reducing sum of base_reflectance at 1550 and 780?
            d_optimal[layer]=d_sweep[layer]
            base_reflectance = (abs(r_init_1550)**2+abs(r_init_780)**2)
    d_sweep[layer]=d_optimal[layer]

wls=[] # Array of wavelengths
T=[] # Array of T values
R=[] # Array of R values  
R_opt=[] # Array of optimized R values
       
for wl in range(430,1700,2):
    n=refractive_index_array(wl)
    r=PS(n,d,wl)               # reflectance of initial values
    r_opt=PS(n,d_optimal,wl)   # reflectance of optimized values
    wls.append(wl)
    R.append(abs(r)**2)
    R_opt.append(abs(r_opt)**2)
#     T.append(1-abs(r)**2)
   
print(d_optimal)
export_array = np.column_stack((wls, R, R_opt))
np.savetxt("tmm.csv", export_array, delimiter=",")

plt.grid(True)    
plt.plot(wls, R, label="Initial")
plt.plot(wls, R_opt, label="optimized")
plt.xlim(430, 1700)
plt.ylim(0.00,1)
plt.xlabel('Wavelength (nm)')
plt.ylabel('Reflectance R')
plt.legend(loc="upper right")
\end{lstlisting}

\section{Fitting code}
\begin{lstlisting}[language=Python, breaklines=true, breakatwhitespace=true]
import numpy as np
import pandas as pd
import matplotlib.pyplot as plt
from scipy.optimize import curve_fit


def tmm_3layer_reflectance(wavelengths, n1_fn, n2_fn, n3_fn, t1, t2, t3, n0=1.0, n_sub=1.5):
    """Calculates reflectance spectrum R(lambda) for a 3-layer thin film stack."""
    wl = np.asarray(wavelengths, dtype=float)
    
    n0_arr = n0(wl) if callable(n0) else np.full_like(wl, n0, dtype=complex)
    nsub_arr = n_sub(wl) if callable(n_sub) else np.full_like(wl, n_sub, dtype=complex)
    
    n_funcs = [n1_fn, n2_fn, n3_fn]
    thicknesses = [t1, t2, t3]
    
    M = np.zeros((len(wl), 2, 2), dtype=complex)
    M[:, 0, 0] = 1.0
    M[:, 1, 1] = 1.0
    
    for n_fn, t in zip(n_funcs, thicknesses):
        n_layer = n_fn(wl) if callable(n_fn) else np.full_like(wl, n_fn, dtype=complex)
        delta = 2.0 * np.pi * n_layer * t / wl
        
        cos_d = np.cos(delta)
        sin_d = np.sin(delta)
        
        m_layer = np.zeros((len(wl), 2, 2), dtype=complex)
        m_layer[:, 0, 0] = cos_d
        m_layer[:, 0, 1] = -1j / n_layer * sin_d
        m_layer[:, 1, 0] = -1j * n_layer * sin_d
        m_layer[:, 1, 1] = cos_d
        
        M = M @ m_layer
        
    B = M[:, 0, 0] + M[:, 0, 1] * nsub_arr
    C = M[:, 1, 0] + M[:, 1, 1] * nsub_arr
    
    r = (n0_arr * B - C) / (n0_arr * B + C)
    return np.abs(r)**2


def fit_n2_parameters(csv_filepath, n1_fn, t1, t2, t3, n0=1.0, n_sub=1.5, p0=[2.0, 5000.0]):
    """
    Fits parameters a and b for n2(wl) = a + b / wl**2 from CSV data.
    Assumes layer 3 refractive index equals layer 1 (n3 = n1).
    
    Parameters:
    - csv_filepath : path to CSV file containing columns 'wavelength' and 'reflectance'
    - n1_fn        : fixed refractive index function or constant for layer 1 (and layer 3)
    - t1, t2, t3   : layer thicknesses (nm)
    - n0           : ambient refractive index
    - n_sub        : substrate refractive index
    - p0           : initial guess [a_init, b_init] for layer 2
    """
    # 1. Load experimental CSV data
    df = pd.read_csv(csv_filepath)
    wl_data = df["wavelength"].values
    R_data = df["reflectance"].values
    
    # Layer 3 has the same refractive index as Layer 1
    n3_fn = n1_fn

    # 2. Define wrapper model for curve_fit (n2 is the unknown function)
    def model_reflectance(wl, a, b):
        n2_fn = lambda w: a + b / (w**2)
        return tmm_3layer_reflectance(wl, n1_fn, n2_fn, n3_fn, t1, t2, t3, n0=n0, n_sub=n_sub)

    # 3. Perform optimization with physical bounds (a > 1, b >= 0)
    bounds = ([1.0, 0.0], [4.0, 1e6])
    popt, pcov = curve_fit(model_reflectance, wl_data, R_data, p0=p0, bounds=bounds)
    
    a_fit, b_fit = popt
    a_err, b_err = np.sqrt(np.diag(pcov))
    
    print("--- Fit Results for Layer 2 ---")
    print(f"a = {a_fit:.5f} +/- {a_err:.5f}")
    print(f"b = {b_fit:.2f} +/- {b_err:.2f} nm^2")
    
    # 4. Plot fitted curve vs experimental data
    R_fit = model_reflectance(wl_data, a_fit, b_fit)
    
    plt.figure(figsize=(8, 5))
    plt.plot(wl_data, R_data, "o", label="CSV Data", color="crimson", alpha=0.6, markersize=4)
    
    # Raw string with .format() avoids LaTeX backslash errors in f-strings
    label_text = r"Fit ($n_2 = {:.3f} + \frac{{ {:.1f} }}{{\lambda^2}}$)".format(a_fit, b_fit)
    plt.plot(wl_data, R_fit, "-", label=label_text, color="black", lw=2)
    
    plt.xlabel("Wavelength (nm)")
    plt.ylabel("Reflectance")
    plt.title("TMM Parameter Fitting for Layer 2 ($n_3 = n_1$)")
    plt.grid(True, linestyle="--", alpha=0.5)
    plt.legend()
    plt.tight_layout()
    plt.show()

    return a_fit, b_fit


# --- Usage Example ---
if __name__ == "__main__":
    sample_wl = np.linspace(400, 800, 50)
    
    # Known layer 1 function (n3 will automatically copy this)
    n1_known = lambda w: np.sqrt(20201856809.3-22607182010.85*w**2/ (w**2-0.1063836985797) +2405325203.746*w**2/ (w**2-0.9998803009815))
    
    # Target values for layer 2: a = 2.1, b = 12000
    n2_target = lambda w: 2.60 + 300000.0 / (w**2)
    
    # Known substrate function
    n_si = lambda w: 4.0333E-12*w**4 - 1.7984E-08*w**3 + 2.9993E-05*w**2 - 2.2437E-02*w + 9.9702
    
    # Generate synthetic experimental data (n1 = n3 = n1_known)
    R_synthetic = tmm_3layer_reflectance(
        sample_wl, n1_known, n2_target, n1_known, t1=208, t2=19, t3=48, n0=1.0, n_sub=n_si
    )
    R_synthetic += np.random.normal(0, 0.003, size=len(sample_wl))  # Add noise
    
    # Save simulated CSV
#     pd.DataFrame({"wavelength": sample_wl, "reflectance": R_synthetic}).to_csv("measured_data.csv", index=False)
    
    # Run parameter fitting for layer 2
    a_opt, b_opt = fit_n2_parameters(
        csv_filepath="measured_data2.csv",
        n1_fn=n1_known,
        t1=208,
        t2=19.0,
        t3=48,
        n0=1.0,
        n_sub=n_si,
        p0=[2.6, 300000.0]  # Initial guess
    )
\end{lstlisting}

% \chapter{Data Processing}
% Data was processed before being added to this document.

\end{appendices}